\chapter{De la distribución de Dirichlet al proceso jerárquico}
\label{app:dirichlet}

\section{Sobre el proceso de Dirichlet y los modelos jerárquicos}

El HDP-HMM relaciona varias distribuciones de transición sin fijar de antemano el número de estados latentes. Para entender esa construcción conviene partir de un problema más sencillo: representar probabilidades sobre un conjunto finito de símbolos y expresar la incertidumbre acerca de ellas. El primer paso utiliza una distribución de Dirichlet; cuando existen varios contextos, sus distribuciones se vinculan mediante una jerarquía. El proceso de Dirichlet extiende esta idea sin prefijar el número de categorías latentes \cite{murphy2012,ferguson1973,teh2006}.

Sea $\mathcal V=\{1,\dots,K\}$ un conjunto finito. Una variable aleatoria $Y$ con valores en $\mathcal V$ sigue una distribución categórica si existe un vector de probabilidades
\[
\mathbf{p}=(p_1,\dots,p_K),\qquad p_k\ge 0,\qquad \sum_{k=1}^K p_k=1,
\]
tal que
\[
P(Y=k\mid \mathbf{p})=p_k.
\]
El vector $\mathbf{p}$ pertenece al simplex de dimensión $K-1$ y contiene el peso asignado a cada categoría posible. Al observar muchas realizaciones para inferir $\mathbf{p}$, los conteos indican qué categorías aparecieron con mayor frecuencia, pero no expresan por sí solos la incertidumbre de la estimación.

La distribución de Dirichlet cumple esa función. Se escribe
\[
\mathbf{p}\sim \mathrm{Dirichlet}(a_1,\dots,a_K),
\]
y define una distribución de probabilidad sobre el conjunto de todos los vectores $\mathbf{p}$ admisibles. Los parámetros $(a_1,\dots,a_K)$ controlan la concentración de masa en el simplex: si todos son grandes, la distribución se concentra más; si son pequeños, permite configuraciones más extremas. En términos intuitivos, la distribución de Dirichlet actúa como una distribución a priori sobre probabilidades categóricas.

Su importancia práctica proviene de la conjugación. Si se observa una muestra de variables categóricas con conteos
\[
(n_1,\dots,n_K),
\]
entonces la distribución posterior de $\mathbf{p}$ vuelve a ser Dirichlet:
\[
\mathbf{p}\mid \text{datos}\sim
\mathrm{Dirichlet}(a_1+n_1,\dots,a_K+n_K).
\]
La actualización bayesiana tiene entonces una forma directa. Cada conteo observado corrige el peso previo asociado a su categoría, y el resultado sigue perteneciendo a la misma familia de distribuciones. La Dirichlet puede entenderse así como una forma de regularizar probabilidades discretas sin perder su interpretación probabilística.

Una sola distribución deja de bastar cuando cada contexto de una secuencia tiene su propio vector de probabilidades. Estimar estos vectores de manera independiente produce estimaciones inestables en los contextos con pocos datos y desaprovecha las similitudes entre las distribuciones locales.

Un modelo jerárquico introduce entonces un nivel superior que gobierna los parámetros locales, tratados a su vez como variables aleatorias. Las distribuciones dejan así de estimarse como piezas aisladas: pueden apartarse unas de otras, pero permanecen vinculadas por una organización global.

En el caso finito, esta idea puede escribirse así. Se introduce primero una distribución global
\[
\mathbf{p}^{(0)}\sim \mathrm{Dirichlet}(b_1,\dots,b_K),
\]
y después se definen distribuciones locales
\[
\mathbf{p}^{(j)}\mid \mathbf{p}^{(0)}
\sim
\mathrm{Dirichlet}(\lambda p^{(0)}_1,\dots,\lambda p^{(0)}_K),
\qquad j=1,\dots,m.
\]
Aquí $\mathbf{p}^{(0)}$ fija un perfil global, mientras que cada $\mathbf{p}^{(j)}$ representa una distribución local que puede apartarse de él, aunque no de manera arbitraria. El parámetro $\lambda$ regula el grado de similitud con el nivel global. Cuando $\lambda$ es grande, las distribuciones locales tienden a parecerse más entre sí; cuando es pequeño, tienen mayor libertad para divergir.

En un HMM hay una distribución inicial, varios renglones de una matriz de transición y, según la formulación, varias distribuciones de emisión. Cada una describe un comportamiento local, aunque todas remiten al mismo repertorio latente. Una estructura jerárquica permite compartir información entre ellas y estabilizar las estimaciones cuando hay pocos datos. En la literatura bayesiana este compromiso se conoce como agrupamiento parcial: los contextos no se estiman de manera independiente, pero tampoco se obliga a que compartan una sola distribución \cite{murphy2012}.

La construcción anterior todavía supone que el número de categorías $K$ es finito y conocido. La distribución de Dirichlet representa la incertidumbre sobre un vector cuya dimensión ya está fijada; si el número de componentes latentes no se conoce de antemano, hace falta pasar de una distribución sobre vectores finitos a una distribución sobre medidas de probabilidad.

El proceso de Dirichlet da este paso \cite{ferguson1973}. En lugar de asignar una distribución a priori a un vector finito $\mathbf{p}$, la asigna a una medida aleatoria $G$. Se mantiene la intuición del caso finito: todavía se representa incertidumbre sobre pesos probabilísticos, pero el repertorio de átomos potenciales ya no tiene que fijarse mediante un entero $K$. El espacio latente puede incorporar tantos componentes como requieran los datos.

Aunque la medida base no sea discreta, las realizaciones del proceso de Dirichlet lo son casi seguramente. Por ello, $G$ puede entenderse como una medida de probabilidad formada por una familia numerable de átomos con pesos aleatorios. La representación de \emph{stick-breaking} introducida por Sethuraman hace explícita esta construcción \cite{sethuraman1994}: la masa total se reparte sucesivamente en una colección numerable de porciones y produce una medida discreta aleatoria.

La jerarquía se conserva en el caso no paramétrico: una medida global sirve de referencia para varias medidas locales, del mismo modo que la distribución global organizaba las distribuciones locales en el caso finito. Con ello se obtiene una familia de distribuciones de transición relacionadas que comparten un repertorio potencialmente infinito de estados latentes \cite{teh2006}.

Cuando no se conoce de antemano la complejidad latente, la jerarquía bayesiana hace que varias distribuciones compartan un repertorio sin imponerles un único comportamiento local. Esta construcción da lugar al HDP-HMM y a su aproximación truncada.
